\begin{frame}{Literature Review (Paper 3)}
\begin{columns}[T]
  \begin{column}{0.25\textwidth}
    \begin{block}{Title of the paper \cite{memoryarch26}}
      \tiny \textbf{Multi-Agent Memory from a Computer Architecture Perspective: Visions and Challenges Ahead}\\[2pt] Zhongming Yu, Naicheng Yu, Hejia Zhang, Wentao Ni, Mingrui Yin, Jiaying Yang, Yujie Zhao, Jishen Zhao (2026).\\[2pt] Problem: multi-agent memory is becoming a systems/architecture bottleneck involving hierarchy, access and consistency.
    \end{block}
    \vspace{0.12cm}
    \begin{block}{Main contributions}
      \tiny
      \begin{itemize}\setlength\itemsep{1pt}
      \item Frames multi-agent memory as a computer-architecture problem.
\item Distinguishes shared and distributed memory paradigms.
\item Proposes a three-layer I/O--cache--memory hierarchy.
\item Identifies cache-sharing and structured-access-control protocol gaps.
      \end{itemize}
    \end{block}
    \vspace{0.12cm}
    \begin{block}{Inputs and Outputs}
      \tiny Conceptual architectural analysis of multi-agent memory, mapping computer-architecture abstractions to agent memory and identifying open protocol/consistency problems.
    \end{block}
  \end{column}

  \begin{column}{0.4\textwidth}
    \begin{block}{Methodology}
    \tiny
    \begin{center}
    \textbf{Architectural / position framework}
    \end{center}
    \begin{itemize}\setlength\itemsep{1pt}
    \item Compare shared vs. distributed memory.
\item Map I/O, cache and memory layers to agent systems.
\item Examine access protocols and consistency requirements.
\item Identify open research directions.
    \end{itemize}
    \end{block}

    \begin{block}{Tools and Technology used}
      \tiny Computer-architecture abstractions; memory hierarchy; access protocols; consistency mechanisms. The paper is a position paper rather than an empirical model benchmark.
    \end{block}
  \end{column}

  \begin{column}{0.25\textwidth}
    \begin{block}{Summary of Results}
      \tiny The paper argues that \textbf{multi-agent memory consistency} is a central open challenge and that hierarchical memory plus explicit protocols are needed for reliable, scalable systems.
    \end{block}
    \vspace{0.12cm}
    \begin{block}{Major Findings}
      \tiny
      \begin{itemize}\setlength\itemsep{1pt}
      \item Memory should be treated as an architectural subsystem.
\item Consistency is not solved merely by providing shared storage.
\item Protocol design and structured access remain open problems.
      \end{itemize}
    \end{block}
  \end{column}
  \end{columns}
\end{frame}
